\documentclass[12pt]{article}
\usepackage[utf8]{inputenc}
\usepackage[T1]{fontenc}
\usepackage[margin=1in]{geometry}
\usepackage{setspace}
\usepackage{xcolor}
\usepackage[colorlinks=true, urlcolor=blue!60!black, linkcolor=blue!60!black]{hyperref}
\setlength{\parindent}{0em}
\setlength{\parskip}{0.45em}
\renewcommand{\baselinestretch}{1.0}
\begin{document}
\begin{center}
{\Large \textit{Plyler v. Doe}, 457 U.S. 202 (1982)\par}
\vspace{0.3em}
{\normalsize Case Brief --- One-Page Edition\par}
\end{center}

\textbf{Memory Jogger.} No denying free public schooling to undocumented children; the law must further a substantial state goal (\href{https://supreme.justia.com/cases/federal/us/457/202/}{457 U.S. 202}).

\textbf{Facts.} A 1975 Texas law withheld state funds for educating children not ``legally admitted'' into the United States and let local districts deny them free enrollment; in 1977 the Tyler Independent School District began charging them a \$1,000 annual tuition fee. Undocumented children of Mexican origin sued.

\textbf{Procedural History.} District courts found the law neither deterred immigration nor improved schools and enjoined it; the Fifth Circuit affirmed.

\textbf{Votes.} Affirmed 5--4. Brennan wrote for the Court; Marshall, Blackmun, and Powell concurred; Burger, with White, Rehnquist, and O'Connor, dissented.

\textbf{Holding.} Undocumented immigrants are ``persons within its jurisdiction.'' Though not a suspect class, and though education is not a fundamental right, denying their children schooling fails unless it furthers a substantial state goal.

\textbf{Reasoning.} Children ``can affect neither their parents' conduct nor their own status,'' and exclusion risks a permanent ``shadow population''; education sustains democratic institutions and its denial brands children with illiteracy. Tuition was a ``ludicrously ineffectual attempt to stem the tide of illegal immigration,'' and no federal policy supported exclusion. Powell invoked the illegitimacy cases; Blackmun likened total denial to denying the vote. Burger called the policy ``senseless'' but said the Court acted as ``Platonic Guardians.''

\textbf{Examples.} \emph{Same-Side:} proof of status to enroll. \emph{Opposite-Side:} a neutral residency rule (\textit{Martinez v. Bynum}). \emph{Fence-Sitter:} status reporting by schools.

\textbf{Critique.} Critics say \textit{Plyler} fused two non-triggers into ad hoc scrutiny better suited to preemption. \textit{Rodriguez} (funding disparities) and \textit{Kadrmas} (bus fees, stressing \textit{Plyler}'s ``unique circumstances'') mark its limits.

\textbf{Key Quotation.} ``Section 21.031 imposes a lifetime hardship on a discrete class of children not accountable for their disabling status. The stigma of illiteracy will mark them for the rest of their lives.'' --- Justice Brennan, 457 U.S. at 223.

\end{document}
